\section{Nozzle Flow \& Thrust}

\subsection{Introduction}
This chapter details the physics of expansion through the nozzle and the resulting thrust generation. The model accounts for isentropic expansion, shifting equilibrium composition, and nozzle efficiency.

\subsection{Nomenclature}
\begin{description}
    \item[\(\thrust\)] Total thrust [\si{N}]
    \item[\(v_{exit}\)] Exit velocity [\si{m/s}]
    \item[\(P_{exit}\)] Exit pressure [\si{Pa}]
    \item[\(P_{amb}\)] Ambient pressure [\si{Pa}]
    \item[\(A_{exit}\)] Exit area [\si{m^2}]
    \item[\(C_f\)] Thrust coefficient [dimensionless]
    \item[\(\Isp\)] Specific impulse [\si{s}]
    \item[\(\epsilon\)] Expansion ratio (\(A_e/A_t\)) [dimensionless]
\end{description}

\subsection{Thrust Equation}
The total thrust is the sum of momentum and pressure components:
\begin{equation}
\thrust = \mdot v_{exit} + (P_{exit} - P_{amb}) A_{exit}
\end{equation}
\coderef{engine/core/nozzle.py}{163, 575}

\subsection{Thrust Coefficient (\(C_f\))}
The thrust can also be expressed using the thrust coefficient:
\begin{equation}
\thrust = C_f P_c A_t
\end{equation}
where \(C_f\) is calculated as:
\begin{equation}
C_f = \frac{\thrust}{P_c A_t}
\end{equation}
\coderef{engine/core/nozzle.py}{166, 614}

\subsection{Exit Conditions}

\subsubsection{Exit Mach Number}
The exit Mach number \(M_{exit}\) is determined from the expansion ratio \(\epsilon\) using the area-Mach relation (supersonic solution):
\begin{equation}
\epsilon = \frac{1}{M_{exit}} \left[ \frac{2}{\gamma+1} \left( 1 + \frac{\gamma-1}{2} M_{exit}^2 \right) \right]^{\frac{\gamma+1}{2(\gamma-1)}}
\end{equation}
\coderef{engine/pipeline/nozzle_dynamics.py}{65}

\subsubsection{Exit Pressure and Temperature}
Isentropic relations relate exit properties to chamber stagnation conditions:
\begin{equation}
P_{exit} = P_c \left( 1 + \frac{\gamma-1}{2} M_{exit}^2 \right)^{-\frac{\gamma}{\gamma-1}}
\end{equation}
\begin{equation}
T_{exit} = T_c \left( 1 + \frac{\gamma-1}{2} M_{exit}^2 \right)^{-1}
\end{equation}
\coderef{engine/core/nozzle.py}{317, 335}

\subsubsection{Exit Velocity}
The ideal exit velocity is:
\begin{equation}
v_{exit,ideal} = M_{exit} \sqrt{\gamma R T_{exit}}
\end{equation}
Applying nozzle efficiency \(\eta_{nozzle}\):
\begin{equation}
v_{exit} = \eta_{nozzle} v_{exit,ideal}
\end{equation}
\coderef{engine/pipeline/nozzle_dynamics.py}{80-84}

\subsection{Specific Impulse (\(\Isp\))}
The specific impulse is defined as:
\begin{equation}
\Isp = \frac{\thrust}{\mdot g_0}
\end{equation}
where \(g_0 = \SI{9.80665}{m/s^2}\).
\coderef{engine/core/nozzle.py}{675}

\subsection{Shifting Equilibrium}
If enabled, the composition shifts during expansion, meaning \(\gamma\) and \(R\) vary. This is solved iteratively to find self-consistent exit properties.
\coderef{engine/core/nozzle.py}{432-500}

\subsection{Bartz Heat Flux Correlation}
The heat flux along the nozzle is estimated using the Bartz correlation:
\begin{equation}
q = 0.026 \left( \frac{\mu^{0.2} c_p}{Pr^{0.6}} \right) \left( \frac{P_c^{0.8}}{D_h^{0.2}} \right) T_{gas}^{0.8} \left( 1 + \frac{\gamma-1}{2} M^2 \right)^{-0.6}
\end{equation}
\coderef{engine/pipeline/nozzle_dynamics.py}{206}

